% GENERATED FILE. Edit canonical lessons, then run npm run generate:books.
% canonical-source-hash: fnv1a64:e6dafc71d39bd45d
% canonical-lessons: JA-C106-mon, JA-C106-hayashi, JA-C106-tsubomi, JA-C106-iwa, JA-C106-koishi

\chapter{Gate, Wood, Bud, Rock, Pebble}
\label{ch:ja-gate-wood-bud-rock-pebble}
\hlchaptermodality{106}

\begin{chapteropening}
\textbf{By the end of this chapter:} \emph{I can say a gate, a wood, a bud, a rock and a pebble.}

Complete the last lesson of chapter 106: i can say a gate, a wood, a bud, a rock and a pebble.
\end{chapteropening}

\section[\texorpdfstring{\ja{もん} (mon)}{もん (mon)}]{\ja{もん} (mon) — a gate}

\label{lesson:JA-C106-mon}



\begin{quote}
Before the new one: say the Japanese for a lamp, then the Japanese for light.
\end{quote}

\subsection*{You'll want to know: \ja{もん}}
\textbf{\ja{もん}} — \emph{mon} — \textquotedblleft{}a gate\textquotedblright{}.

\begin{grammarlens}[title={What you've built}]
One more word for this chapter.
\end{grammarlens}

\subsection*{Your turn}
\begin{itemize}
\raggedright
  \item \emph{Say it:} \emph{mon}
  \item \emph{Say it:} \emph{mon}, once more
  \item \emph{Say it:} say \emph{mon}
  \item \emph{Recall:} say the Japanese for a lamp, then the Japanese for light, then say \emph{mon} again
\end{itemize}

\begin{tcolorbox}[breakable,colback=teal!4,colframe=teal!35!black,title={Before you move on}]
What does \ja{もん} mean? (\textquotedblleft{}A gate\textquotedblright{}.) Say it once more.
\end{tcolorbox}
\section[\texorpdfstring{\ja{はやし} (hayashi)}{はやし (hayashi)}]{\ja{はやし} (hayashi) — a wood, a grove}

\label{lesson:JA-C106-hayashi}



\begin{quote}
Before the new one: say the Japanese for light, then the Japanese for a gate.
\end{quote}

\subsection*{You'll want to know: \ja{はやし}}
\textbf{\ja{はやし}} — \emph{hayashi} — \textquotedblleft{}a wood, a grove\textquotedblright{}.

\begin{grammarlens}[title={What you've built}]
One more word for this chapter.
\end{grammarlens}

\subsection*{Your turn}
\begin{itemize}
\raggedright
  \item \emph{Say it:} \emph{hayashi}
  \item \emph{Say it:} \emph{hayashi}, once more
  \item \emph{Say it:} say \emph{hayashi}
  \item \emph{Recall:} say the Japanese for light, then the Japanese for a gate, then say \emph{hayashi} again
\end{itemize}

\begin{tcolorbox}[breakable,colback=teal!4,colframe=teal!35!black,title={Before you move on}]
What does \ja{はやし} mean? (\textquotedblleft{}A wood, a grove\textquotedblright{}.) Say it once more.
\end{tcolorbox}
\section[\texorpdfstring{\ja{つぼみ} (tsubomi)}{つぼみ (tsubomi)}]{\ja{つぼみ} (tsubomi) — a bud}

\label{lesson:JA-C106-tsubomi}



\begin{quote}
Before the new one: say the Japanese for a gate, then the Japanese for a wood.
\end{quote}

\subsection*{You'll want to know: \ja{つぼみ}}
\textbf{\ja{つぼみ}} — \emph{tsubomi} — \textquotedblleft{}a bud\textquotedblright{}.

\begin{grammarlens}[title={What you've built}]
One more word for this chapter.
\end{grammarlens}

\subsection*{Your turn}
\begin{itemize}
\raggedright
  \item \emph{Say it:} \emph{tsubomi}
  \item \emph{Say it:} \emph{tsubomi}, once more
  \item \emph{Say it:} say \emph{tsubomi}
  \item \emph{Recall:} say the Japanese for a gate, then the Japanese for a wood, then say \emph{tsubomi} again
\end{itemize}

\begin{tcolorbox}[breakable,colback=teal!4,colframe=teal!35!black,title={Before you move on}]
What does \ja{つぼみ} mean? (\textquotedblleft{}A bud\textquotedblright{}.) Say it once more.
\end{tcolorbox}
\section[\texorpdfstring{\ja{いわ} (iwa)}{いわ (iwa)}]{\ja{いわ} (iwa) — a rock}

\label{lesson:JA-C106-iwa}



\begin{quote}
Before the new one: say the Japanese for a wood, then the Japanese for a bud.
\end{quote}

\subsection*{You'll want to know: \ja{いわ}}
\textbf{\ja{いわ}} — \emph{iwa} — \textquotedblleft{}a rock\textquotedblright{}.

\begin{grammarlens}[title={What you've built}]
One more word for this chapter.
\end{grammarlens}

\subsection*{Your turn}
\begin{itemize}
\raggedright
  \item \emph{Say it:} \emph{iwa}
  \item \emph{Say it:} \emph{iwa}, once more
  \item \emph{Say it:} say \emph{iwa}
  \item \emph{Recall:} say the Japanese for a wood, then the Japanese for a bud, then say \emph{iwa} again
\end{itemize}

\begin{tcolorbox}[breakable,colback=teal!4,colframe=teal!35!black,title={Before you move on}]
What does \ja{いわ} mean? (\textquotedblleft{}A rock\textquotedblright{}.) Say it once more.
\end{tcolorbox}
\section[\texorpdfstring{\ja{こいし} (koishi)}{こいし (koishi)}]{\ja{こいし} (koishi) — a pebble}

\label{lesson:JA-C106-koishi}



\begin{quote}
Before the new one: say the Japanese for a bud, then the Japanese for a rock.
\end{quote}

\subsection*{You'll want to know: \ja{こいし}}
\textbf{\ja{こいし}} — \emph{koishi} — \textquotedblleft{}a pebble\textquotedblright{}.

\begin{grammarlens}[title={What you've built}]
One more word for this chapter.
\end{grammarlens}

\subsection*{Your turn}
\begin{itemize}
\raggedright
  \item \emph{Say it:} \emph{koishi}
  \item \emph{Say it:} \emph{koishi}, once more
  \item \emph{Say it:} say \emph{koishi}
  \item \emph{Recall:} say the Japanese for a bud, then the Japanese for a rock, then say \emph{koishi} again
\end{itemize}

\begin{tcolorbox}[breakable,colback=teal!4,colframe=teal!35!black,title={Before you move on}]
What does \ja{こいし} mean? (\textquotedblleft{}A pebble\textquotedblright{}.) Say it once more.
\end{tcolorbox}
